\begin{tikzpicture}[
    x=1cm,
    y=1.35cm,
    >=stealth,
    font=\small\fontspec{Times New Roman}\songti,
    card/.style={draw, rounded corners=4pt, line width=0.45pt, fill=white},
    chip/.style={
        draw,
        rounded corners=2pt,
        line width=0.35pt,
        fill=white,
        align=center,
        inner xsep=5pt,
        inner ysep=5pt
    },
    bluecard/.style={card, draw=blue!45!black, fill=blue!4},
    cyancard/.style={card, draw=cyan!50!black, fill=cyan!5},
    orangecard/.style={card, draw=orange!65!black, fill=orange!6},
    purplecard/.style={card, draw=purple!45!black, fill=purple!5},
    redcard/.style={card, draw=red!45!black, fill=red!5},
    panel/.style={draw=black!15, rounded corners=5pt, line width=0.35pt, fill=gray!1},
    arrow/.style={->, line width=0.55pt, draw=black!70},
    note/.style={font=\scriptsize\fontspec{Times New Roman}\songti, align=center}
]

\draw[panel] (-0.25,0.00) rectangle (21.25,9.00);
\node[font=\large\fontspec{Times New Roman}\songti] at (10.50,8.55)
    {CBTPU-DeepJDOT 多证据置信均衡伪标签控制机制};

% 基础输入
\draw[bluecard] (0.45,2.15) rectangle (3.25,7.75);
\node at (1.85,7.43) {基础信息};
\node[chip, minimum width=2.20cm] at (1.85,6.05)
    {TPU-DeepJDOT\\[-1pt]$\gamma^\ast,\mu_c$};
\node[chip, minimum width=2.20cm] at (1.85,4.60)
    {目标域样本\\[-1pt]$x_j^t$};
\node[note, text width=2.30cm] at (1.85,3.25)
    {继承第三章传输与原型结构};

% 教师-学生增强
\draw[cyancard] (3.95,2.15) rectangle (7.35,7.75);
\node at (5.65,7.43) {教师--学生增强};
\node[chip, minimum width=2.65cm] at (5.65,6.15)
    {弱增强 + 教师锚点\\[-1pt]$a_w(x_j^t),\bar{\Theta}$};
\node[chip, minimum width=2.65cm] at (5.65,4.60)
    {强增强 + 学生\\[-1pt]$a_s(x_j^t),\Theta$};
\node[note, text width=2.70cm] at (5.65,3.25)
    {加载 TPU 检查点；冻结时不参与 EMA};

% 三路证据
\draw[orangecard] (8.05,2.15) rectangle (11.15,7.75);
\node at (9.60,7.43) {三路语义证据};
\node[chip, minimum width=2.35cm] at (9.60,6.20)
    {传输语义\\[-1pt]$q^{ot}$};
\node[chip, minimum width=2.35cm] at (9.60,4.95)
    {教师语义\\[-1pt]$q^{cls}$};
\node[chip, minimum width=2.35cm] at (9.60,3.70)
    {原型语义\\[-1pt]$q^{proto}$};

% 接收门控
\draw[purplecard] (11.85,2.15) rectangle (15.45,7.75);
\node at (13.65,7.43) {可靠接收门控};
\node[chip, minimum width=2.75cm] at (13.65,6.20)
    {多证据融合\\[-1pt]$q^{mix}$};
\node[chip, minimum width=2.75cm, text width=2.65cm] at (13.65,4.90)
    {接收集合 $\mathcal{A}$\\[-1pt]一致 / JS / $\vartheta(\nu)$ / 熵};
\node[chip, minimum width=2.75cm] at (13.65,3.55)
    {接收比例上限\\[-1pt]$\kappa,\,H(q^{ot})$};
\node[note, text width=2.90cm] at (13.65,2.60)
    {只让高可靠目标样本进入伪监督};

% 均衡目标学习
\draw[redcard] (16.15,2.15) rectangle (20.85,7.75);
\node at (18.50,7.43) {均衡目标学习};
\node[chip, minimum width=3.65cm] at (18.50,6.20)
    {类别均衡伪监督\\[-1pt]$\mathcal{L}_{pl},\ \hat{\pi}_c$};
\node[chip, minimum width=3.65cm] at (18.50,4.95)
    {强弱增强一致性\\[-1pt]$\mathcal{L}_{cons}$};
\node[chip, minimum width=3.65cm] at (18.50,3.70)
    {信息最大化辅助\\[-1pt]$\mathcal{L}_{im}$};
\node[note, text width=3.75cm] at (18.50,2.60)
    {降低伪标签噪声与类别坍缩};

% 目标域学习
\draw[redcard] (3.95,0.50) rectangle (20.85,1.35);
\node at (12.40,0.92)
    {$\mathcal{L}_{CBTPU}
    =\mathcal{L}_{TPU}
    +\lambda_{pl}\mathcal{L}_{pl}
    +\lambda_{cons}\mathcal{L}_{cons}
    +\lambda_{im}\mathcal{L}_{im}$};

\draw[arrow] (3.25,4.95) -- (3.95,4.95);
\draw[arrow] (7.35,4.95) -- (8.05,4.95);
\draw[arrow] (11.15,4.95) -- (11.85,4.95);
\draw[arrow] (15.45,4.95) -- (16.15,4.95);
\draw[arrow] (18.50,2.15) -- (18.50,1.35);

\end{tikzpicture}
